\documentclass[a4paper,12pt]{article}
\usepackage{../sty/config}
\usepackage{amsmath}
\usepackage{amsfonts}
\usepackage{amssymb}

\begin{document}

% Configuración de encabezado manual para que coincida con el estilo
\pagestyle{fancy}

\begin{center}
    \Huge \textbf{Actividad Asincrónica}\\
    \Large Ley de Avogadro y Mezclas de Gases
    \vspace{1cm}
\end{center}

\section*{Instrucciones}
\begin{itemize}
    \item Esta actividad es de carácter asincrónico. 
    \item \textbf{Tiempo estimado:} 60 minutos.
    \item Utilice la constante de los gases ideales $R = 0,0821 \frac{\text{L atm}}{\text{mol K}}$ o su equivalente según las unidades requeridas.
    \item Se recomienda mostrar todo el procedimiento de resolución.
\end{itemize}

\section*{Actividades}

\begin{enumerate}[label=\textbf{\arabic*.}, leftmargin=*]

    \item \textbf{Fundamentos Teóricos (10 min)}\\
    Explique la diferencia conceptual entre la \textbf{Ley de Dalton} de las presiones parciales y la \textbf{Ley de Amagat} de los volúmenes parciales. ¿En qué condiciones de un gas ideal ambas leyes resultan consistentes entre sí?

    \item \textbf{Cálculo de Estado (10 min)}\\
    Un recipiente de $15\text{ L}$ contiene un gas a una presión de $2,5\text{ atm}$ y una temperatura de $27^\circ\text{C}$. Calcule el número de moles de gas presentes en el recipiente.

    \item \textbf{Relación Masa-Presión (10 min)}\\
    Se tiene una muestra de $80\text{ g}$ de un gas cuya masa molar es de $40\text{ g/mol}$. Si el gas ocupa un volumen de $12\text{ L}$ a una temperatura de $300\text{ K}$, ¿cuál es la presión que ejerce el gas?

    \item \textbf{Mezclas y Presiones Parciales (10 min)}\\
    Una mezcla gaseosa está compuesta por $0,5\text{ moles}$ de Helio ($He$) y $1,5\text{ moles}$ de Argón ($Ar$). Si la presión total de la mezcla es de $4\text{ atm}$, determine la presión parcial de cada uno de los gases.

    \item \textbf{Transformaciones Isobáricas (10 min)}\\
    Un tanque de gas de $5\text{ L}$ se encuentra a una temperatura de $25^\circ\text{C}$. Si el tanque se calienta hasta alcanzar los $80^\circ\text{C}$ manteniendo la presión constante, ¿cuál será el nuevo volumen que ocuparía esa misma cantidad de gas?

    \item \textbf{Mezcla de Gases y Volumen Total (10 min)}\\
    Una mezcla contiene $32\text{ kg}$ de Oxígeno ($O_2$, $M=32\text{ g/mol}$) y $28\text{ kg}$ de Nitrógeno ($N_2$, $M=28\text{ g/mol}$). Si la mezcla se encuentra a una temperatura de $300\text{ K}$ y una presión de $2\text{ atm}$, calcule el volumen total que ocupa la mezcla.

\end{enumerate}

\vfill
\begin{center}
    \textit{Fin de la actividad. ¡Éxitos!}
\end{center}

\end{document}
